% Part: first-order-logic
% Chapter: sequent-calculus
% Section: provability-quantifiers

\documentclass[../../../include/open-logic-section]{subfiles}

\begin{document}

\olfileid[pt]{fol}{seq}{qpr}
\olsection{\usetoken{S}{derivability} e os Quantificadores}

\begin{explain}
  O teorema da completude também exige que as regras do cálculo de
  sequentes produzam os fatos sobre~$\Proves$ estabelecidos nesta
  seção.
\end{explain}

\begin{thm}
\ollabel{thm:strong-generalization} Se $c$ é uma constante que não ocorre
em $\Gamma$ ou $!A(x)$ e $\Gamma \Proves !A(c)$, então $\Gamma
\Proves \lforall[x][!A(x)]$.
\end{thm}

\begin{proof}
Seja $\pi_0$ uma $\Log{LK}$-!!{derivation} de $\Gamma_0 \Sequent !A(c)$
para algum $\Gamma_0 \subseteq \Gamma$ finito.  Acrescentando uma
inferência $\RightR{\lforall}$, obtemos !!a{derivation} de $\Gamma_0 \Sequent
\lforall[x][!A(x)]$, já que $c$ não ocorre em $\Gamma$ ou $!A(x)$
e, assim, a condição de variável própria é satisfeita.
\end{proof}

\begin{prop}
\ollabel{prop:provability-quantifiers}
\begin{tagenumerate}{prvEx,prvAll}
\tagitem{prvEx}{$!A(t) \Proves \lexists[x][!A(x)]$.}{}
\tagitem{prvAll}{$\lforall[x][!A(x)] \Proves !A(t)$.}{}
\end{tagenumerate}
\end{prop}

\begin{proof}
\begin{tagenumerate}{prvEx,prvAll}
\tagitem{prvEx}{O sequente $!A(t) \Sequent \lexists[x][!A(x)]$ é
  !!{derivable}:
  \begin{prooftree}
    \Axiom$!A(t) \fCenter !A(t)$
    \RightLabel{\RightR{\lexists}}
    \UnaryInf$!A(t) \fCenter \lexists[x][!A(x)]$
    \end{prooftree}}{}
\tagitem{prvAll}{O sequente $\lforall[x][!A(x)] \Sequent !A(t)$ é
  !!{derivable}:
  \begin{prooftree}
    \Axiom$!A(t) \fCenter !A(t)$
    \RightLabel{\LeftR{\lforall}}
    \UnaryInf$\lforall[x][!A(x)] \fCenter !A(t)$
    \end{prooftree}}{}
\end{tagenumerate}
\end{proof}

\end{document}
